\subsection{What Should the Model Be Trained Against?}
\label{section:fc_value_aware}

Throughout the preceding subsections the model has been trained against a likelihood or reconstruction loss. The Recurrent State-Space Model line uses a pixel reconstruction term, a reward log-likelihood, and a Kullback-Leibler consistency term between prior and posterior; the original Dyna model is a tabular maximum-likelihood estimate of the empirical transition; the controller-model architecture of \citet{HaSchmidhuber2018} fits its variational autoencoder and its mixture-density recurrent network to observation and next-latent likelihoods respectively. In each case the loss rewards the model for matching observed distributions, regardless of how those distributions enter the value calculation the agent actually optimizes. The objection raised in the present subsection is that the agent does not care about predicting the next observation perfectly; it cares about getting the value function right. This subsection traces the line that took that observation seriously, from value-aware model learning in \citet{FarahmandVAML2017} through the value-equivalence principle of \citet{GrimmVALEQ2020,GrimmPVE2021}, the value-targeted regret bound of \citet{AyoubVTR2020}, the MuZero instantiation of \citet{Schrittwieser2020MuZero}, and the recent calibration analysis of \citet{VoelckerCalib2025}.

\subsubsection{Value-aware model learning}
\label{section:fc_value_aware_vaml}

\citet{FarahmandVAML2017} formalize the objection as a loss. Fix a value function class $\mathcal{V}$, a data distribution $\mu$ over state-action pairs, and a model class within which the learner selects an estimate $\widehat P$ of the true transition kernel $P^\star$. The value-aware model learning loss is
\[ \mathcal{L}_{\text{VAML}}(\widehat P) = \sup_{V \in \mathcal{V}} \Big\Vert (\widehat P V)(\cdot) - (P^\star V)(\cdot) \Big\Vert_\mu, \]
where $(\widehat P V)(s, a) = \int V(s') \widehat P(\mathrm{d} s' \mid s, a)$ is the expected next-state value under the learned model and $(P^\star V)(s, a)$ is the same quantity under the true kernel. The loss penalizes the worst-case discrepancy in expected next-state value across $\mathcal{V}$, measured under $\mu$. A model that places mass on features of the next observation irrelevant to every value function in $\mathcal{V}$ is not penalized; a model that distorts the value expectation is. \citet{FarahmandIterVAML2018} embed the same idea inside approximate value iteration, replacing the supremum over a value class with the realized value iterate at each step, which yields the practical IterVAML estimator and a fitted error bound that propagates through the Bellman backups rather than through the model alone. \citet{AsadiWasserstein2018} read the construction through optimal transport, observing that under a Lipschitz assumption on the value class the VAML loss reduces to a Wasserstein distance between $\widehat P$ and $P^\star$, which connects the value-aware loss to the substantial empirical literature on Wasserstein generative modeling.

VAML is the formal statement of the intuition that the model should be wrong in ways that do not matter for the value functional. This is the technical answer to the conceptual openness of Sutton's original Dyna architecture, in which the model was permitted to be probabilistic and oftentimes incorrect (\S\ref{section:fc_dyna_q_architecture}) without specifying what kind of error was tolerable. The likelihood-trained Recurrent State-Space Model line (\S\ref{section:fc_rssm}) commits to a particular notion of error, namely the reconstruction loss on pixels and the Kullback-Leibler divergence on latents, and violates the value-aware principle whenever the pixel distribution is rich in detail orthogonal to the reward. Whether this matters in practice depends on the task; the empirical claim of the value-aware line is that it does, and that the gap widens as observation dimensionality grows.

\subsubsection{Value equivalence and MuZero}
\label{section:fc_value_aware_ve}

\citet{GrimmVALEQ2020} sharpen the VAML idea into a clean equivalence relation. Fix a policy set $\Pi$ and a value function set $\mathcal{V}$. Two transition kernels $\widehat P$ and $P^\star$ are value-equivalent with respect to $(\Pi, \mathcal{V})$ when their Bellman operators agree on every pair $(\pi, V) \in \Pi \times \mathcal{V}$, that is when $(T^\pi_{\widehat P} V)(s) = (T^\pi_{P^\star} V)(s)$ for all $s$ and all $(\pi, V)$ in the chosen sets. The equivalence class shrinks monotonically as $\Pi$ and $\mathcal{V}$ grow, collapsing to the singleton $\{P^\star\}$ in the limit. The practical implication is that a small model class can be sufficient for a restricted set of policies and value functions, even if it would be insufficient as a generative model of the environment. \citet{GrimmPVE2021} subsequently introduce proper value equivalence, which restricts the value class to fixed points of the chosen Bellman operators and gives the variant of value equivalence that aligns most directly with the MuZero training objective. \citet{AyoubVTR2020} push the value-aware idea into the regret-bound literature with value-targeted regression, which trains the model against value targets and matches the dimension-dependent rates of optimism-based reinforcement learning under standard assumptions.

The empirical instantiation is MuZero. \citet{Schrittwieser2020MuZero} train a latent recurrence jointly with a policy network and a value network against three losses applied at the next step of the recurrence, namely a reward prediction loss, a value prediction loss, and a policy prediction loss matched to the output of a Monte Carlo tree search rooted at the same latent. There is no reconstruction loss; the latent dynamics are never asked to decode observations. The transition kernel itself has no explicit target, since the only quantities that enter any loss are the reward, the value, and the policy at the rollout step. Empirically the construction matches AlphaZero on Go, chess, and shogi without access to the rules of those games, and reaches state-of-the-art performance on the Atari benchmark. \citet{Antonoglou2022StochMuZero} extend the construction to stochastic transitions by inserting a learned afterstate node, the game-tree analogue of a chance node, between the agent action and the next state. Read in economics terms, value-aware model learning treats the world model as a task-relevant misspecification, kept correct on the value functional even when it is wrong on the full observation distribution.\footnote{\citet{VoelckerCalib2025} revisit this family and observe that the standard sampled losses, including the canonical MuZero objective, are uncalibrated surrogate losses whose minima do not in general coincide with the minima of the underlying target loss on the value function; value-equivalence is a property of an in-distribution policy and value class, and a model trained purely against value targets has no incentive to remain accurate once the policy drifts. The paper proposes a calibrated variant that restores the correctness property in expectation and leaves the broader question of value-aware accuracy under policy drift open.}\footnote{The same move appears in the operations-research line on decision-focused learning, in which a predictor is trained with a loss aligned to the downstream decision quality rather than to a statistical scoring rule on the forecast itself \citep{DontiAmosKolter2017}. The inner solver there is a convex program rather than a Bellman backup, but the conceptual move, to align the loss with what the decision-maker actually needs and to treat likelihood as the wrong default whenever the model class is misspecified, is shared.}
